\subsection{THETA3-C21: split CM formal factors and higher-dimensional theta ratios}\label{C21}

\paragraph{Statement.}
Let $A$ be a $g$-dimensional abelian variety over a number field,
with CM by an order $\mathcal O$ in a CM field $E$ of degree
$2g$.  Write $F$ for the maximal totally real subfield of $E$.
Let $p\geq5$ split completely in $F$, and suppose every prime
$\nu$ of $F$ above $p$ is unramified and inert in $E$.
Assume that $\mathcal O$ is maximal at $p$.
After a finite extension choose a good abelian scheme
$\mathcal A/R$, with $R$ the integers of a finite extension
of $\Q_p$, and define all CM endomorphisms over it.
Enlarge the field further to contain the CM embeddings whenever
needed.  Then
\begin{equation}\label{C21-formal-product}
 \widehat{\mathcal A}\simeq
       \prod_{\nu\mid p}G_\nu,
\end{equation}
where every $G_\nu$ is an integral one-dimensional formal
module for the unramified quadratic extension of $\Q_p$,
with height-two special fiber.

Choose integral invariant differentials on the factors and
write $D_\nu$ for their dual invariant derivations, normalized
by integral parameters with $D_\nu t_\nu(0)=1$.
For $f\in R[[t_\nu:\nu\mid p]]$ and a multi-index
$\mathbf n=(n_\nu)$, put
\[
 b_p(n)=\max\left\{0,
            \left\lfloor\frac{np}{p^2-1}\right\rfloor-1\right\}.
\]
Then the intrinsic mixed-derivative estimate is
\begin{equation}\label{C21-mixed-bound}
 v_p\left(\left(\prod_{\nu\mid p}D_\nu^{n_\nu}\right)f(0)\right)
       \geq\sum_{\nu\mid p}b_p(n_\nu).
\end{equation}
In particular the right side is at least
$p|\mathbf n|/(p^2-1)-2r$, where $r$ is the number of
nonzero entries in $\mathbf n$.

There is a direct theta application with a concrete integral
hypothesis.  Suppose $s_1,\ldots,s_m$ and
$u_1,\ldots,u_m$ are integral sections of the same power
$\mathscr L^{\otimes k}$ of a line bundle on $\mathcal A$,
and every $u_i$ is a unit at the identity in an integral local
frame.  Their equal-degree ratio
\begin{equation}\label{C21-section-ratio}
 f=\frac{\prod_i s_i}{\prod_i u_i}
\end{equation}
belongs to the integral formal ring and satisfies
\eqref{C21-mixed-bound}.  The same conclusion holds after
an explicitly chosen common prime-to-$p$ isogeny pullback
has trivialized the flat torsion twists of translated theta
sections, provided the resulting sections have the stated
integral normalization and the denominator is a unit at the
identity.  No assertion that all higher-dimensional tame
theta values are automatically units is used.

\paragraph{Integral CM idempotents and their heights.}
The local splitting hypotheses give the exact algebra
\[
 \mathcal O\otimes\Z_p
       =\mathcal O_E\otimes\Z_p
       \simeq\prod_{\nu\mid p}\mathcal O_{E_\nu},
 \qquad E_\nu/\Q_p\text{ unramified quadratic}.
\]
Let $e_\nu$ be its mutually orthogonal idempotents.
These act on $\mathcal A[p^n]$ for every $n$: the action
factors through $\mathcal O/p^n\mathcal O$, and the
idempotents are compatible as $n$ varies.
An idempotent on a finite flat commutative group splits it
into its image and its complementary image, with inverse
maps given by the idempotent and its complement.
Thus
\[
 \mathcal A[p^\infty]=\prod_{\nu\mid p}\mathcal G_\nu,
       \qquad \mathcal G_\nu=e_\nu\mathcal A[p^\infty].
\]
The factors are finite flat at every level; equivalently,
their coordinate algebras are direct retracts of the finite
flat coordinate algebra before splitting.

The rational Betti homology of $A$ is a rank-one $E$-module:
it is a faithful $E$-vector space of rational dimension $2g$.
Using complex uniformization, its $p$-adic completion identifies
the rational Tate module as a rank-one
$E\otimes\Q_p$-module.  Hence every $\mathcal G_\nu$
has height two.
The CM type selects one embedding out of every conjugate
pair of embeddings of $E$.  Because $p$ splits completely
in $F$, the two embeddings extending a fixed $\nu$ form
exactly such a pair.  Therefore
$e_\nu\operatorname{Lie}(\mathcal A)$ is an integral direct
summand of rank one.  Each $\mathcal G_\nu$ has dimension
one, and its $\mathcal O_{E_\nu}$ action on the tangent
line is through one specified embedding into $R$.

Here is why no \'etale part remains in its special fiber.
The Tate module of that \'etale part is a free $\Z_p$-module
of some rank $h_{\rm et}\leq2$, stable under
$\mathcal O_{E_\nu}$.  Its rationalization is an
$E_\nu$-vector space, so $h_{\rm et}$ is even.
It cannot be two: the remaining connected part would have
height zero, making the whole factor \'etale of dimension
zero, whereas its dimension is one.  Thus $h_{\rm et}=0$.
The special fiber is connected.  Equivalently its geometric
$p$-torsion, an $\mathbb F_{p^2}$-module, has only the
identity point.  All the finite kernels are therefore supported
at the identity in the special fiber.

\paragraph{A direct construction of the formal product.}
The completed local CM action exists on
$\widehat{\mathcal A}$ without an algebraization claim for
the idempotents.  In arbitrary integral coordinates
$t_1,\ldots,t_g$, put $I=(p,t_1,\ldots,t_g)$.
Each component of $[p](\mathbf t)$ equals $pt_i$ plus
terms of total degree at least two.  The same induction as
in C19 gives $[p^n](\mathbf t)\in I^{n+1}$ componentwise.
Approximating local CM elements by global ones consequently
converges $I$-adically to formal endomorphisms, respecting
addition, composition and tangent maps.

Choose the initial coordinates so that their tangent
covectors respect the rank-one direct-sum decomposition
defined by the $e_\nu$.  Set
\[
 X_\nu=t_\nu\circ e_\nu.
\]
The linear part of the tuple $(X_\nu)$ is the identity.
The integral formal inverse-function theorem therefore makes
it a system of integral coordinates.  Orthogonality gives
$X_\mu(e_\nu P)=0$ for $\mu\ne\nu$, and
$X_\nu(e_\nu P)=X_\nu(P)$.
Thus the image of $e_\nu$ is precisely the coordinate axis
on which all other $X_\mu$ vanish.  Because the $e_\nu$ are
group homomorphisms and their sum is the identity, the group
law is the product of its restrictions to these axes.
This proves \eqref{C21-formal-product} with each factor
smooth of dimension one.

The connectedness proved above identifies the $p^n$-kernels
of these formal factors with the entire finite flat groups
$e_\nu\mathcal A[p^n]$.  In particular $[p]$ on each
$G_\nu$ is finite flat of degree $p^2$, with height-two
special fiber.  Identifying $\mathcal O_{E_\nu}$ through
its tangent embedding makes its action strict, exactly as
required by the formal-module generalization of C19.
The exact basic references are Tate, \emph{$p$-Divisible
Groups}, \S1.4, p.~160 (connected--\'etale sequence),
\S2.2, Proposition~1, p.~162 with proof pp.~162--164
(connected groups and divisible formal groups), the
dimension definition following (4), p.~164, and
\S2.3, Proposition~3 and Example~(b), p.~166
(invariance of height and dimension under reduction and
the formal completion of an abelian scheme).
The idempotent splitting and coordinate construction have
been given explicitly above.

\paragraph{Mixed derivatives.}
C19 gives the full one-variable operator divisibility
\[
 D_\nu^nR[[X_\nu]]\subset
              p^{b_p(n)}R[[X_\nu]].
\]
Apply it coefficientwise in the other formal variables.
The factors are independent, so their derivations commute,
and successive application multiplies these integral powers
of $p$.  This proves \eqref{C21-mixed-bound}.
For $n>0$, $b_p(n)\geq np/(p^2-1)-2$; summing over the
nonzero entries gives the displayed total-degree consequence.
As in C19, finite-jet approximation extends the assertion to
integral $\C_p$ coefficients on the fixed base-changed group.
Rescaling any differential by an integral unit preserves the
valuation bound.  A nonunit analytic period rescaling must
instead be retained explicitly in the derivative formula.

\paragraph{Integral sections, common frames, and torsion twists.}
Trivialize $\mathscr L$ on the formal neighborhood of the
identity by an integral frame.  The numerator sections in
\eqref{C21-section-ratio} then have integral coefficients,
and every denominator series has a unit constant term.
The quotient is consequently an integral formal function.
Changing the common frame multiplies each section by the
same factor to the same tensor degree, so that factor cancels
exactly.  In the canonical complex theta trivialization this
includes the common quadratic exponential potential.
Thus the canonical theta ratio represents the same formal
function in the exponential coordinates normalized by the
chosen invariant differentials.  The analytic matching is
the degree-zero case of A8.  If the supplied theta sections
are rescaled by constants, their ratio is correspondingly
rescaled by their product quotient $c$; the resulting bound
is \eqref{C21-mixed-bound} plus $v_p(c)$.

For translated sections, $\tau_T^*\mathscr L$ differs from
$\mathscr L$ by the flat line bundle attached to
$\phi_{\mathscr L}(T)\in\mathcal A^\vee$.
If all the translations have prime-to-$p$ torsion order,
choose a common prime-to-$p$ integer $M$ killing these
dual torsion points.  Pullback by $[M]$ trivializes their
flat twists: on rigidified degree-zero line bundles it
corresponds to multiplication by $M$ on the dual variety.
Choose these integral identifications and retain the
integrality and unit-denominator conditions in the statement.
The pulled-back equal-degree ratio is then a rational function
on the isogeny cover whose completion is integral.
This assertion does not require it to descend to the
unpulled-back abelian variety.

On the formal neighborhood $[M]$ is an integral isomorphism,
with inverse having integral coefficients because its tangent
matrix is $M$ times the identity and $M\in R^\times$.
Also
\[
 D_\nu[M]^*=M[M]^*D_\nu.
\]
For a mixed derivative at the identity, passage to the
pullback therefore multiplies its value by
$M^{|\mathbf n|}$, a unit.  The same valuation estimate
holds in the original local translated-theta coordinates
with the specified identifications. \hfill$\square$

\paragraph{Remaining scope.}
This supplies the strong local estimate on an actual
higher-dimensional CM formal group at the stated split-in-$F$,
inert-in-$E$ primes, and on canonical theta ratios meeting
the explicit integral-section hypothesis.  Positive tensor
degree retains the canonical connection; C22 shows why
that issue cannot be removed by declaring every theta
section to have the ratio estimate.  Neither the formal
idempotents nor this analytic estimate construct the
special-cycle correspondence required in C17.
